\section{浸入与嵌入}\label{sec:immersions-embeddings}

本节研究光滑流形之间的可微映射 $f:M\to N$ 的局部行为. 我们已经看到, 若 $\dim M=\dim N$, 且 $f$ 将 $p$ 的一个邻域微分同胚地映到 $f(p)$ 的一个邻域, 则称 $f$ 在点 $p\in M$ 处为局部微分同胚. 这时, 它的导数 $(\mathrm{d}f)_p:T_pM\to T_{f(p)}N$ 必为同构 [参见 Exercise~\ref{ex:tangent-space}(\ref{ex:smoothness-and-inverse-of-pushforward})]. 反过来, 若 $(\mathrm{d}f)_p$ 是同构, 则由反函数定理可知 $f$ 是局部微分同胚 (参见第~\RGref{subsec:notes-inverse-function-theorem}~节). 因此, 要检验 $f$ 是否将 $p$ 的一个邻域微分同胚地映到 $f(p)$ 的一个邻域, 只需检验 $(\mathrm{d}f)_p$ 的局部表示的行列式是否非零.

当 $\dim M<\dim N$ 时, 至多只能要求 $(\mathrm{d}f)_p:T_pM\to T_{f(p)}N$ 是单射. 这时称 $f$ \textbf{在 $p$ 处为浸入}. 若 $f$ 在 $M$ 的每一点处均为浸入, 则称它为\textbf{浸入}. 局部而言, 每个浸入在相差微分同胚的意义下, 都是 $\mathbb{R}^m$ 到 $\mathbb{R}^n$ ($m<n$) 的\textbf{典范浸入}, 即将点 $(x^1,\ldots,x^m)$ 映到 $(x^1,\ldots,x^m,0,\ldots,0)$. 这一结果称为\textbf{局部浸入定理}.

\begin{RiemannTheorem}\label{thm:local-immersion-normal-form}
设 $f:M\to N$ 在 $p\in M$ 处为浸入. 则在 $p$ 和 $f(p)$ 附近存在局部坐标, 使得 $f$ 在这些坐标下为典范浸入.
\end{RiemannTheorem}

\begin{RiemannProof}
设 $(U,\varphi)$ 和 $(V,\psi)$ 是 $p$ 和 $q=f(p)$ 附近的参数化. 为简单起见, 假设 $\varphi(0)=p$ 且 $\psi(0)=q$. 由于 $f$ 是浸入, $(\mathrm{d}\widehat f)_0:\mathbb{R}^m\to\mathbb{R}^n$ 是单射 (其中 $\widehat f:=\psi^{-1}\circ f\circ\varphi$ 是 $f$ 的局部坐标表达式). 因此, 可以假定 (必要时更换 $\mathbb{R}^n$ 的基) 这个线性变换由如下 $n\times m$ 矩阵表示:
\[
  \left(\begin{array}{c} I_{m\times m}\\ \hdashline 0\end{array}\right),
\]
其中 $I_{m\times m}$ 是 $m\times m$ 恒等矩阵. 因此, 映射
\begin{align*}
  F:U\times\mathbb{R}^{n-m}&\to\mathbb{R}^n\\
  (x^1,\ldots,x^n)&\mapsto
  \widehat f(x^1,\ldots,x^m)+(0,\ldots,0,x^{m+1},\ldots,x^n),
\end{align*}
的导数 $(\mathrm{d}F)_0:\mathbb{R}^n\to\mathbb{R}^n$ 由下列矩阵给出:
\[
  \left(\begin{array}{c:c}
  I_{m\times m}&0\\ \hdashline
  0&I_{(n-m)\times(n-m)}
  \end{array}\right)=I_{n\times n}.
\]
应用反函数定理可知, $F$ 在 $0$ 处是局部微分同胚. 这意味着 $\psi\circ F$ 在 $0$ 处也是局部微分同胚, 因而 $\psi\circ F$ 是 $N$ 在 $q$ 附近的另一个参数化. 将 $\mathbb{R}^m$ 到 $\mathbb{R}^n$ 的典范浸入记为 $j$, 则 $\widehat f=F\circ j\Leftrightarrow f=\psi\circ F\circ j\circ\varphi^{-1}$, 因而下图交换:
\[
  \begin{tikzcd}[ampersand replacement=\&]
    M\supset\varphi(\widetilde U) \arrow[r,"f"] \& (\psi\circ F)(\widetilde V)\subset N\\
    \mathbb{R}^m\supset\widetilde U \arrow[u,"\varphi"] \arrow[r,"j"] \&
    \widetilde V\subset\mathbb{R}^n \arrow[u,"\psi\circ F"']
  \end{tikzcd}
\]
(其中 $\widetilde U\subset U$ 和 $\widetilde V\subset V$ 是可能需要缩小的开集). 因此, $f$ 在这些新坐标下是典范浸入.
\end{RiemannProof}

\begin{RiemannRemark}\label{rem:immersion-on-neighborhood}
由局部浸入定理可知, 在点 $p\in M$ 处为浸入的映射, 在 $p$ 的一个邻域上也是浸入.
\end{RiemannRemark}

若浸入 $f:M\to N$ 又是从 $M$ 到赋有子空间拓扑的像 $f(M)\subset N$ 的同胚, 则称它为\textbf{嵌入}. 我们把证明局部浸入定理蕴含每个浸入在局部都是嵌入留作习题.

\RGfigure[.55\textwidth]{injective-immersion-not-embedding}{不是嵌入的单射浸入}

\begin{RiemannExample}\label{ex:immersions-and-embeddings}
\begin{enumerate}[label=(\arabic*),ref=\arabic*,leftmargin=*]
\item 由 $f(t)=(t^2,t^3)$ 给出的映射 $f:\mathbb{R}\to\mathbb{R}^2$ 在 $t=0$ 处不是浸入.
\item\label{ex:periodic-plane-immersion} 由 $f(t)=(\cos t,\sin 2t)$ 定义的映射 $f:\mathbb{R}\to\mathbb{R}^2$ 是浸入, 但不是嵌入 (它不是单射).
\item 设 $g:\mathbb{R}\to\mathbb{R}$ 为函数 $g(t)=2\arctan(t)+\pi/2$. 若 $f$ 为 Example~\ref{ex:immersions-and-embeddings}(\ref{ex:periodic-plane-immersion}) 中定义的映射, 则 $h:=f\circ g$ 是单射浸入, 但不是嵌入. 实际上, Fig.~\ref{fig:injective-immersion-not-embedding} 中的集合 $S=h(\mathbb{R})$ 并不是 $\mathbb{R}$ 到 $\mathbb{R}^2$ 的某个嵌入的像. 图中的箭头表示曲线在原点处与自身任意接近, 但从不自交. 若在 $\mathbb{R}$ 和 $\mathbb{R}^2$ 上取通常拓扑, 则 $\mathbb{R}$ 中包含 $0$ 的有界开集的像, 对子空间拓扑而言并不是 $h(\mathbb{R})$ 中的开集, 所以 $h^{-1}$ 不连续.
\item 由 $f(t)=(e^t\cos t,e^t\sin t)$ 给出的映射 $f:\mathbb{R}\to\mathbb{R}^2$ 是 $\mathbb{R}$ 到 $\mathbb{R}^2$ 的嵌入.
\end{enumerate}
\end{RiemannExample}

若 $M\subset N$ 且包含映射 $i:M\hookrightarrow N$ 是嵌入, 就称 $M$ 为 $N$ 的\textbf{子流形}. 因此, 嵌入 $f:M\to N$ 将 $M$ 微分同胚地映到 $N$ 的一个子流形. $f(M)$ 上的坐标图就是在 $N$ 上适当选取的坐标图到 $f(M)$ 的限制 [参见 Exercise~\ref{ex:immersions-embeddings}(\ref{ex:coordinate-submanifold-characterization})].

若可微映射 $f:M\to N$ 的 $(\mathrm{d}f)_p$ 是满射, 则称 $f$ \textbf{在 $p$ 处为淹没}. 注意, 这时必有 $m\geq n$. 若 $f$ 在 $M$ 的每一点处均为淹没, 则称它为\textbf{淹没}. 局部而言, 每个淹没都是 $\mathbb{R}^m$ 到前 $n$ 个因子的标准投影.

\begin{RiemannTheorem}\label{thm:local-submersion-normal-form}
设 $f:M\to N$ 在 $p\in M$ 处为淹没. 则在 $p$ 和 $f(p)$ 附近存在局部坐标, 使得 $f$ 在这些坐标下为标准投影.
\end{RiemannTheorem}

\begin{RiemannProof}
考虑 $p$ 和 $f(p)$ 附近的参数化 $(U,\varphi)$ 和 $(V,\psi)$, 使得 $f(\varphi(U))\subset\psi(V)$, $\varphi(0)=p$ 且 $\psi(0)=f(p)$. $f$ 在局部坐标下由 $\widehat f:=\psi^{-1}\circ f\circ\varphi$ 给出. 由于 $(\mathrm{d}f)_p$ 是满射, $(\mathrm{d}\widehat f)_0:\mathbb{R}^m\to\mathbb{R}^n$ 是满射线性变换. 对 $\mathbb{R}^n$ 作线性坐标变换, 可以假定 $(\mathrm{d}\widehat f)_0=(I_{n\times n}\mid *)$. 如同局部浸入定理的证明, 我们引入一个辅助映射 $F$, 以便应用反函数定理:
\begin{align*}
  F:U\subset\mathbb{R}^m&\to\mathbb{R}^m\\
  (x^1,\ldots,x^m)&\mapsto(\widehat f(x^1,\ldots,x^m),x^{n+1},\ldots,x^m).
\end{align*}
它在 $0$ 处的导数是由下式给出的线性映射:
\[
  (\mathrm{d}F)_0=\left(\begin{array}{c:c}
  I_{n\times n}&*\\ \hdashline
  0&I_{(m-n)\times(m-n)}
  \end{array}\right).
\]
由反函数定理可知, $F$ 在 $0$ 处是局部微分同胚, 即它将这个点的某个开邻域 $\widetilde U\subset U$ 微分同胚地映到 $\mathbb{R}^m$ 中包含 $0$ 的一个开集 $W$. 若 $\pi_1:\mathbb{R}^m\to\mathbb{R}^n$ 是到前 $n$ 个因子的标准投影, 则 $\pi_1\circ F=\widehat f$, 因此
\[
  \widehat f\circ F^{-1}=\pi_1:W\to\mathbb{R}^n.
\]
于是, 用 $\widetilde\varphi:=\varphi\circ F^{-1}$ 替代 $\varphi$, 得到使 $f$ 成为到前 $n$ 个因子的标准投影 $\pi_1$ 的坐标:
\[
  \psi^{-1}\circ f\circ\widetilde\varphi
  =\psi^{-1}\circ f\circ\varphi\circ F^{-1}
  =\widehat f\circ F^{-1}=\pi_1.
\]
\end{RiemannProof}

\begin{RiemannRemark}\label{rem:rank-theorem}
这一结果通常与局部浸入定理一同陈述, 合称为\textbf{秩定理} (例如见 \RGcite{Boo03}).
\end{RiemannRemark}

设 $f:M\to N$ 是维数分别为 $m$ 和 $n$ 的光滑流形之间的可微映射. 若 $(\mathrm{d}f)_p$ 是满射, 则称点 $p\in M$ 是 $f$ 的\textbf{正则点}. 若 $f^{-1}(q)$ 中的每一点都是正则点, 则称点 $q\in N$ 是 $f$ 的\textbf{正则值}. 不是正则点的点 $p\in M$ 称为 $f$ 的\textbf{临界点}, 相应的值 $f(p)$ 称为\textbf{临界值}. 注意, 若 $f$ 存在正则值, 则 $m\geq n$. 取正则值的原像可以得到可微流形.

\begin{RiemannTheorem}\label{thm:submanifold-of-regular-level-set}
设 $q\in N$ 是 $f:M\to N$ 的正则值, 并假定水平集 $L:=f^{-1}(q)=\{p\in M\mid f(p)=q\}$ 非空. 则 $L$ 是 $M$ 的子流形, 且对所有 $p\in L$, 有 $T_pL=\ker(\mathrm{d}f)_p\subset T_pM$.
\end{RiemannTheorem}

\begin{RiemannProof}
对每个点 $p\in f^{-1}(q)$, 在 $p$ 和 $q$ 附近选取参数化 $(U,\varphi)$ 和 $(V,\psi)$, 使得 $f$ 是到前 $n$ 个因子的标准投影 $\pi_1$, 且 $\varphi(0)=p$, $\psi(0)=q$ (参见 Theorem~\ref{thm:local-submersion-normal-form}). 接着按以下方式为 $L:=f^{-1}(q)$ 构造可微结构: 从 $M$ 的这些参数化中分别取出集合 $U$. 由于 $f\circ\varphi=\psi\circ\pi_1$, 有
\begin{align*}
  \varphi^{-1}\bigl(f^{-1}(q)\bigr)
    &=\pi_1^{-1}\bigl(\psi^{-1}(q)\bigr)=\pi_1^{-1}(0)\\
    &=\{(0,\ldots,0,x^{n+1},\ldots,x^m)\mid x^{n+1},\ldots,x^m\in\mathbb{R}\},
\end{align*}
于是
\[
  \widetilde U:=\varphi^{-1}(L)
  =\{(x^1,\ldots,x^m)\in U\mid x^1=\cdots=x^n=0\};
\]
因此, 若取 $\pi_2:\mathbb{R}^m\to\mathbb{R}^{m-n}$ 为到后 $m-n$ 个因子的标准投影, 取 $j:\mathbb{R}^{m-n}\to\mathbb{R}^m$ 为由
\[
  j(x^1,\ldots,x^{m-n})=(0,\ldots,0,x^1,\ldots,x^{m-n}),
\]
给出的浸入, 则族 $\{(\pi_2(\widetilde U),\varphi\circ j)\}$ 是 $L$ 的一个图册.

此外, 包含映射 $i:L\to M$ 是嵌入. 实际上, 若 $A$ 是 $L$ 中包含在一个坐标邻域内的开集, 则
\[
  A=\varphi\left(\left(\mathbb{R}^n\times(\varphi\circ j)^{-1}(A)\right)\cap U\right)\cap L
\]
是 $L$ 的子空间拓扑下的开集.

现在证明 $T_pL=\ker(\mathrm{d}f)_p$. 为此, 对每个 $v\in T_pL$, 考虑 $L$ 上满足 $c(0)=p$ 且 $\dot c(0)=v$ 的曲线 $c$. 于是对每个 $t$, 有 $(f\circ c)(t)=q$, 因此
\[
  \frac{\mathrm{d}}{\mathrm{d}t}(f\circ c)(0)=0
  \Leftrightarrow(\mathrm{d}f)_p\dot c(0)=(\mathrm{d}f)_pv=0,
\]
故 $v\in\ker(\mathrm{d}f)_p$. 由于 $\dim T_pL=\dim\bigl(\ker(\mathrm{d}f)_p\bigr)=m-n$, 结论得证.
\end{RiemannProof}

给定一个可微流形, 可以问它是否能嵌入到某个 $\mathbb{R}^K$ ($K\in\mathbb{N}$) 中. 下面由 Whitney 在 [\RGref{bib:Whi44a}, \RGref{bib:Whi44b}] 中证明的定理回答了这一问题, 称为 \textbf{Whitney 嵌入定理}.

\begin{RiemannTheorem}[Whitney]\label{thm:whitney-embedding}
任何 $n$ 维光滑流形 $M$ 都可以嵌入 $\mathbb{R}^{2n}$ (且当 $n>1$ 时, 可以浸入 $\mathbb{R}^{2n-1}$).\hfill\openbox
\end{RiemannTheorem}

\begin{RiemannRemark}\label{rem:manifold-as-euclidean-submanifold}
由 Whitney 嵌入定理可知, 任何 $n$ 维光滑流形 $M$ 都与 $\mathbb{R}^{2n}$ 的一个子流形微分同胚.
\end{RiemannRemark}

\begin{RiemannExercise}\label{ex:immersions-embeddings}
\begin{enumerate}[label=(\arabic*),ref=\arabic*,leftmargin=*]
\item\label{ex:parametrizations-are-embeddings} 证明任何参数化 $\varphi:U\subset\mathbb{R}^m\to M$ 都是 $U$ 到 $M$ 的嵌入.
\item\label{ex:immersion-is-locally-an-embedding} 证明每个浸入在局部都是嵌入. 即若 $f:M\to N$ 是浸入且 $p\in M$, 则存在包含 $p$ 的开集 $W\subset M$, 使得 $f|_W$ 是嵌入.
\item\label{ex:coordinate-submanifold-characterization} 设 $N$ 是流形. 证明: $M\subset N$ 是 $N$ 的一个 $m$ 维子流形, 当且仅当对每个 $p\in M$, 在 $N$ 上 $p$ 的附近都有一个坐标系 $x:W\to\mathbb{R}^n$, 使得 $M\cap W$ 由方程 $x^{m+1}=\cdots=x^n=0$ 定义.
\item\label{ex:sphere-as-regular-level-set} 考虑球面
\[
  S^n=\{x\in\mathbb{R}^{n+1}\mid(x^1)^2+\cdots+(x^{n+1})^2=1\}.
\]
证明 $S^n$ 是 $\mathbb{R}^{n+1}$ 的一个 $n$ 维子流形, 且
\[
  T_xS^n=\{v\in\mathbb{R}^{n+1}\mid\langle x,v\rangle=0\},
\]
其中 $\langle\cdot,\cdot\rangle$ 是 $\mathbb{R}^n$ 上的通常内积.
\item\label{ex:restriction-to-submanifolds} 设 $f:M\to N$ 是光滑流形之间的可微映射, 并考虑子流形 $V\subset M$ 和 $W\subset N$. 证明若 $f(V)\subset W$, 则 $f:V\to W$ 也是可微映射.
\item\label{ex:compact-injective-immersion} 设 $f:M\to N$ 是单射浸入. 证明若 $M$ 紧致, 则 $f(M)$ 是 $N$ 的子流形.
\end{enumerate}
\end{RiemannExercise}
